\section{Determinism and Freedom}

The rule is fully deterministic. There are no dice anywhere in the base, and at first glance that looks like a
cage. It is not. Being determined is not the same as being predestined, and freedom was never the opposite of
being determined in the first place. Freedom is the opposite of being coerced and the opposite of being random.
Between coercion and randomness sits a real third thing, a choice that is fully caused, authored from within, and
unknowable until it is lived. This chapter locates that third thing precisely. It is the compatibilist position,
and the strongest argument against the deterministic base is named plainly along the way.

\subsection{Determined yet not predestined}

Two facts pry the two notions apart. The first is that there is no finished total state. The crystal grows without
end, with new vibes born at the frontier every beat, and whether or not there was a first cell is left open. No
single moment holds all the information of all of time. A future that exists nowhere in advance cannot be looked
up, cannot be read off, and cannot be the script a puppet walks. The second fact is computational
irreducibility. For a universal system there is in general no closed form for the far future faster than running it
beat by beat \cite{margenstern2013}. The only computer that can compute the future is the system itself, running
forward.

\begin{principle}[Determined is not predestined]
The future is caused but neither pre-written nor forecastable. With no finished total state and no shortcut faster
than running the system, the future is determined as it happens rather than sitting ready in advance. Apparent
randomness is the local view, since every event is fixed by the whole mesh it sits in while no local part holds
enough of the mesh to predict it.
\end{principle}

The phrase to hold onto is determined by the whole and opaque to the part. Every event is fixed, but it is fixed
by far more of the structure than any local observer can hold, which is why prediction fails even though
causation does not.

\subsection{Determined yet self-authored}

A self is a self-maintaining pattern. The rule is reversible and has no attractors, so a self persists by doing
work, which is the same reason memory costs work. A choice is that self resolving an urge, a pull that arrives
through perception from a deeper and slower part of the mesh, into one definite next act. Four things are true of
the resolving at once, and together they form the definition of freedom used here.

\begin{itemize}
  \item \emph{It is determined.} The same self under the same urge resolves the same way. That is what makes the
  choice yours rather than noise. A choice that came out differently each time would be a malfunction, not
  authorship.
  \item \emph{It is yours.} A different self, given the same urge, resolves it differently. This is genuine
  downward causation, with the macro-pattern shaping which micro-tones come next.
  \item \emph{A deeper self has more say.} A richer, more integrated self dominates the urge rather than being
  shoved by it. Agency scales with integration.
  \item \emph{It is irreducible.} No one, not even the universe, can know how a self will resolve without letting
  it resolve. There is no oracle that skips the computation.
\end{itemize}

\begin{principle}[Free means caused by you]
Freedom is not the absence of cause. It is causation that runs through the self's own integrated structure. When
the self models and steers itself, the causal loop closes and the self is caused by the self, which is what the
words I decided name, a real loop in the dynamics rather than an exception to them.
\end{principle}

The compatibilist core is that free does not mean uncaused, it means caused by you. The puppet's strings are real,
but they are attached to the self's own integrated structure rather than to anything outside it. The loop is the
freedom.

\subsection{The strongest objection, and the micro door}

The deterministic base has a serious adversary, and it should be stated rather than skirted. The free will theorem
of Conway and Kochen argues that, under modest assumptions, if experimenters have a kind of freedom then so do the
particles, in the sense that a particle's response is not a function of the prior information available anywhere in
its past light cone \cite{conwaykochen2009}. Taken at face value this is in tension with a fully deterministic
microdynamics. The fair move is to concede the micro door. The theory does not claim to settle the status of
freedom at the level of single elementary events, and it does not need to. Its claim about freedom is pitched at
the macro level, the level of integrated selves, where the compatibilist account above does its work. A self's
freedom is the integrated pattern computing its own resolution, and that account survives whatever the right
reading of the micro theorem turns out to be.

\subsection{The will, finite and spendable}

The will is a directed pump, a sustained push that biases the self's resolution against the local gradient. It is
not an exception to the rule but an input the rule respects, and its effects are measured \cite{pollard2026vibetest}.
With the will active a self moves toward a chosen target and away from a threat. At a fork between a small reward
now and a larger reward later, a self with enough will delays gratification. The pump can be depleted, and when it
runs out the self relapses to the greedy choice at a sharp threshold rather than fading gradually. So the will is a
finite, spendable, directed capacity, and that finiteness is the texture of real willpower as it is actually
experienced.

\begin{proposition}[The will is a finite directed pump]
The will is a sustained directed push that biases resolution against the local gradient, respected by the rule
rather than excepted from it. It enables movement toward targets and away from threats and enables delay of
gratification, and it depletes, with relapse to the greedy choice at a sharp threshold once it is spent.
\end{proposition}


A deterministic base does not abolish freedom. It locates it. Freedom is the integrated self computing its own
resolution, unknowable in advance even to the universe, and genuinely yours, one path but authored, surprising even
to its author. Freedom is not a sixth base thing. It is what the four base things do when a self is integrated
enough to fold back on itself and pump against its own gradient. The landscape-of-valleys image often used to
picture the weighing of options is a way to visualize the choice, not a claim about the reversible microdynamics,
and the no-finished-state argument is kept compatible with the open question of whether the substrate is finite or
infinite. That last question stays open. What does not depend on it is the central point, that determination at the
base and freedom at the level of selves are not in conflict.
